%%=================================================================
%% CHAPTER 3: Dark Electromagnetic Duality
%% Dark Energy as the T-Breaking Component of Dark Matter
%% Source: dark-energy-T-breaking/ (within Secluded-Majorana-SIDM)
%%=================================================================

\section{The Central Analogy}
\label{ch3:analogy}

In classical electromagnetism, a static charge produces only an electric field
$\vec{E}$. Motion reveals a second component — the magnetic field $\vec{B}$ —
which is not a new entity but a different projection of the same tensor $F_{\mu\nu}$.
The critical property: $\vec{B}$ is T-odd while $\vec{E}$ is T-even.

We propose the dark sector obeys the same structure:

\begin{table}[h!]
\centering
\caption{Dark Electromagnetic Duality correspondence.}
\begin{tabular}{ll}
\toprule
Electromagnetism & Dark Sector \\
\midrule
$\vec{E}$ (T-even) & $y_s = y\cos(\sigma/f)$ — attraction, SIDM clustering \\
$\vec{B}$ (T-odd) & $y_p = y\sin(\sigma/f)$ — repulsion, DE acceleration \\
$F_{\mu\nu}$ & $\mathcal{Y} = y_s + iy_p = ye^{i\gamma^5\sigma/f}$ \\
Lorentz boost mixes $E \leftrightarrow B$ & $\sigma$ field rotates $y_s \leftrightarrow y_p$ \\
$\vec{B}$ does no work ($\vec{F}\perp\vec{v}$) & $\gamma^5$ prevents $y_p$ from contributing to mass \\
Like charges repel electrically & $y_s$ causes DM–DM attraction (SIDM) \\
Parallel currents attract magnetically & $y_p$ causes vacuum repulsion (DE) \\
\bottomrule
\end{tabular}
\end{table}

\textbf{Thesis}: Dark energy is the T-violating component of the dark matter
interaction, made manifest through CP violation in the Majorana sector —
exactly as magnetism is the T-violating component of electrostatics,
made manifest through relative motion.

\section{The Extended Lagrangian}
\label{ch3:lagrangian}

Starting from the SIDM Lagrangian (Chapter~1), we promote the CP phase to a
dynamical field $\sigma$ — a dark pseudo-Goldstone boson:
\begin{equation}
  y_s + iy_p \;\longrightarrow\; ye^{i\gamma^5\sigma/f}
  = y\cos\!\left(\frac{\sigma}{f}\right) + iy\sin\!\left(\frac{\sigma}{f}\right)\gamma^5
\end{equation}

The full unified Lagrangian:
\begin{equation}
  \boxed{\Lc = \underbrace{\frac{1}{2}\bar{\chi}(i\slashed{\partial} - \mchi)\chi
    + \frac{1}{2}(\partial\phi)^2 - V_0(\phi)}_{\rm SIDM\;(Ch.1)}
    - \frac{1}{2}\bar{\chi}\!\left[y\cos\!\tfrac{\sigma}{f}
      + iy\sin\!\tfrac{\sigma}{f}\gamma^5\right]\!\chi\,\phi
    + \underbrace{\frac{1}{2}(\partial\sigma)^2 - V_{\rm bare}(\sigma)}_{\rm dark\;axion}}
\end{equation}

Parameters: $y$ fixed by SIDM+relic (Chapter~1); $f \sim 0.2\,M_{\rm Pl}$ from
$V_\sigma \sim \rho_\Lambda$ (see §\ref{ch3:scale}); all other parameters determined.

\section{The Universal Relic Angle}
\label{ch3:thetarel}

Two simultaneous constraints:
\begin{align}
  \text{SIDM (Ch.1):} \quad & \alphass = \alpha \\
  \text{Relic (Ch.1):} \quad & \alphass\alphapp = \frac{\alpha^2}{8}
\end{align}

These fix $\alphapp = \alpha/8$, hence:
\begin{equation}
  \boxed{\thetarel = \arctan\!\sqrt{\frac{\alphapp}{\alphass}}
         = \arctan\!\frac{1}{\sqrt{8}} = 19.47^\circ}
\end{equation}

This angle is \emph{universal} — independent of $(m_\chi, m_\phi, \alpha)$.
It is the unique equilibrium where the same coupling
simultaneously satisfies DM self-interactions and the observed relic abundance.

\section{Scale of the Dark Energy Potential}
\label{ch3:scale}

The effective potential at $\theta = \thetarel$:
\begin{align}
  V_\sigma &\approx V_{CW}(\thetarel) \cdot \left(\frac{f}{M_{\rm Pl}}\right)^4 \\
  V_\sigma &\approx \rho_\Lambda \quad\Leftrightarrow\quad f \approx 0.2\,M_{\rm Pl}
\end{align}

Numerically, for the MAP benchmark:
\begin{equation}
  f_{\rm MAP} = \frac{M_{\rm Pl}}{(V_{CW}/\rho_\Lambda)^{1/4}}
              \approx 6.6\times10^{17}\,{\rm GeV} = 0.27\,M_{\rm Pl}
\end{equation}

The sub-Planckian value $f \sim 0.2$–$0.3\,M_{\rm Pl}$ is natural in string
compactifications and aligns with the weak gravity conjecture.

\section{Fifth Force and Screening}
\label{ch3:fifthforce}

The $\sigma$ field mediates a fifth force with coupling:
\begin{equation}
  \beta = \frac{M_{\rm Pl}}{f} \approx 5.0 \quad \text{(universal for all BPs)}
\end{equation}

Existing Planck+BAO constraints give $\beta < 0.066$ for an unscreened
coupled quintessence field. Resolution requires:
(a) Chameleon screening in dense environments, or
(b) $\sigma$ coupling only to dark matter (secluded),
    not to baryons — consistent with the secluded dark sector hypothesis.

Since $\phi$ has no tree-level SM coupling (§\ref{ch1:model}), $\sigma$ — being
a phase of the $\chi$–$\phi$ system — naturally inherits this seclusion.
Baryonic fifth-force constraints do not directly apply.

\section{Observational Signatures}
\label{ch3:signatures}

\begin{enumerate}
  \item \textbf{Equation of state}: The dark axion potential predicts
        $w_{\rm DE} = -1 + \mathcal{O}(\theta_i^2/3)$ near $\theta_i \sim 2$ rad.
        Current DESI constraint ($w_0 = -0.827 \pm 0.063$) is consistent at
        $\sim 1.6\sigma$ from the model prediction ($w_0 \approx -0.727$).

  \item \textbf{Hubble tension}: $H_0 = 67.4$ km/s/Mpc (from Chapter~2)
        is consistent with CMB measurements and does not worsen the H$_0$ tension
        relative to $\Lambda$CDM.

  \item \textbf{CP violation in SIDM}: The $y_s/y_p$ ratio is in principle
        measurable through indirect detection asymmetries at
        $\mathcal{O}(\sin^2\theta_{relic}) \approx 10\%$ level.

  \item \textbf{Null prediction}: No baryonic fifth force.
        No cross-coupling between $\sigma$ and SM particles at tree level.
\end{enumerate}

\section{Summary: From SIDM to Dark Energy in One Lagrangian}
\label{ch3:summary}

The three parts of this paper are not separate models. They are three
projections of one Lagrangian:

\begin{itemize}
  \item \textbf{T-even projection} ($y_s$, $\phi$): DM self-interactions,
        velocity-dependent cross sections, Fornax cores, relic density.
        This is Chapter~1.
  \item \textbf{Path-integral structure}: VPM = fluctuation determinant,
        CW vacuum, Matsubara freeze-out, dark QCD $\to H_0$.
        This is Chapter~2.
  \item \textbf{T-odd projection} ($y_p$, $\sigma$): vacuum energy, DE equation
        of state, universal relic angle. This is Chapter~3.
\end{itemize}

The unification is not formal — the same coupling $y$, the same mass $m_\chi$,
and the same mediator mass $m_\phi$ that explain the rotation curves of SPARC
galaxies and the survival of Fornax globular clusters also set the scale of
dark energy and predict $H_0 = 67.4$ km/s/Mpc.
